\section{民事诉讼与民事诉讼法}
\begin{enumerate}
  \item 诉讼。两造（双方）平等对抗，法官居中裁判，解决他们之间的纠纷。
  \item 对抗的两造性质不同以及解决的纠纷的性质不同。
        \begin{enumerate}
          \item 民事诉讼。平等主体。人身、财产、权利、义务纠纷。
          \item 刑事诉讼。解决被告刑事责任的有无和大小问题的纠纷。
          \item 行政诉讼。行政相对人和行政机关，行政行为合法性。
        \end{enumerate}
  \item 民事纠纷的解决途径。
        \begin{enumerate}
          \item 和解。只有双方当事人参与，没有任何第三方的介入。私力救济。
          \item 调解。调解组织、调解委员会。\al{法院调解 3 方，和解 2 方 \dui。参与调解的大爷大妈那叫其他诉讼参与人}。
          \item 仲裁。社会组织、群众组织。社会救济。
          \item 诉讼。国家公权力，审判权。公力救济。
        \end{enumerate}
  \item 民事审判程序分为诉讼和非讼。
        \begin{enumerate}
          \item 解决纠纷的叫诉讼。一审、二审、再审。
          \item 不解决纠纷的叫非讼。特别程序、督促程序、公示催告程序、执行程序。
          \item 选民资格。解决的是公民的政治权利，不属于民事审判，既不是诉讼，也不是非讼。
        \end{enumerate}
  \item 民事诉讼法。
        \begin{enumerate}
          \item 整个民事诉讼法规定的程序不都是民事诉讼程序。
          \item 属性。法律体系 \yjt 基本法。社会关系 \yjt 部门法。规范内容 \yjt 程序法。国家审判权 \yjt 公法。
          \item 效力。在中国进行民事诉讼，适用中国的民事诉讼法。国家主权问题。作为公法，行使国家审判权。
        \end{enumerate}
\end{enumerate}
